% BHU PYQ -- Transcribed & Corrected
% Paper: COMMD11: Elementary marketing
% Exam: Under Graduate Multidisciplinary Course Semester I Examination 2024-25 (NEP)
% Department: Faculty of Commerce

\documentclass[11pt,a4paper]{article}
\usepackage[a4paper,top=2.5cm,bottom=2.5cm,left=2.8cm,right=2.8cm,headheight=14pt]{geometry}
\usepackage{amsmath,amssymb}
\usepackage[T1]{fontenc}
\usepackage[utf8]{inputenc}
\usepackage{lmodern,microtype,fancyhdr,enumitem,hyperref,lastpage,parskip}

\hypersetup{colorlinks=true,urlcolor=black,linkcolor=black,pdftitle={COMMD11: Elementary Marketing -- BHU Sem I 2024-25 (NEP MDC)}}
\pagestyle{fancy}\fancyhf{}
\fancyhead[L]{\small\textit{Banaras Hindu University}}
\fancyhead[R]{\small\textit{COMMD11: Elementary marketing}}
\fancyfoot[C]{\small Page \thepage\ of \pageref{LastPage}}

\newlist{parts}{enumerate}{2}
\setlist[parts,1]{label=(\alph*),leftmargin=*,itemsep=4pt,topsep=2pt}
\newcommand{\pts}[1]{\hfill\textit{[#1]}}

\begin{document}

\begin{center}
  {\large\bfseries BANARAS HINDU UNIVERSITY}\\[2pt]
  {\normalsize Under Graduate Multidisciplinary Course Semester I Examination 2024-25 (NEP)}\\[6pt]
  \rule{0.8\linewidth}{0.5pt}\\[6pt]
  {\large\bfseries FACULTY OF COMMERCE}\\[4pt]
  {\large\bfseries Paper Code: COMMD11 --- Elementary marketing}\\[4pt]
  {\normalsize Time: 2:30 Hours\hfill Maximum Marks: 70}
\end{center}

\rule{\linewidth}{0.4pt}\smallskip
\noindent\textit{Instruction: Write your Roll No. at the top immediately on the receipt of this question paper.}\\[2pt]
\noindent\textit{Note: Answer all questions. All questions carry equal marks. / सभी प्रश्नों के उत्तर दीजिए। सभी प्रश्नों के अंक समान हैं।}
\bigskip

\noindent\textbf{Question 1} \\
'Marketing is ascertaining, creating and satisfying the wants of the people for getting profit'. In the light of this statement explain meaning and scope of marketing.\\
'लाभार्जन हेतु लोगों की आवश्यकताओं का पता लगाना, सृजन करना एवं तृप्त करना ही विपणन है।' इस कथन के संदर्भ में विपणन के अर्थ एवं क्षेत्र को समझाइए।

\begin{center}\textbf{OR / अथवा}\end{center}

Define market segmentation with the help of example and evaluate the various bases of market segmentation.\\
बाजार विभक्तीकरण को उदाहरण की मदद से परिभाषित कीजिए और बाजार विभक्तीकरण के विभिन्न आधारों का मूल्यांकन कीजिए।

\medskip\rule{\linewidth}{0.2pt}\medskip

\noindent\textbf{Question 2} \\
What do you mean by product? Examine the factors which are kept in mind while taking product decision.\\
उत्पाद से आपका क्या अभिप्राय है? उन घटकों का परीक्षण कीजिए जिन्हें उत्पाद निर्णयों के समय ध्यान में रखा जाता है।

\begin{center}\textbf{OR / अथवा}\end{center}

Analyse the different stages of product life cycle.\\
उत्पाद चक्र के विभिन्न चरणों का विश्लेषण कीजिए।

\medskip\rule{\linewidth}{0.2pt}\medskip

\noindent\textbf{Question 3} \\
'Without prices, there is no marketing'. In the light of this statement explain the concept and significance of pricing.\\
'विपणन बिना मूल्यों के नहीं होता है।' इस कथन के संदर्भ में मूल्यन की अवधारणा एवं महत्ता को समझाइए।

\begin{center}\textbf{OR / अथवा}\end{center}

Examine the different methods of pricing which are used to determine the price of a product.\\
किसी वस्तु की कीमत निर्धारण के लिए विभिन्न मूल्यन विधियों का परीक्षण कीजिए।

\medskip\rule{\linewidth}{0.2pt}\medskip

\noindent\textbf{Question 4} \\
What is channel of distribution? Discuss different factors which affect the selection of marketing channel.\\
वितरण वाहिका क्या है? विपणन वाहिका के चुनाव को प्रभावित करने वाले विभिन्न घटकों की विवेचना कीजिए।

\begin{center}\textbf{OR / अथवा}\end{center}

What is advertising? Explain its advantages and disadvantages.\\
विज्ञापन क्या होता है? इसके लाभ और हानियों को समझाइए।

\medskip\rule{\linewidth}{0.2pt}\medskip

\noindent\textbf{Question 5} \\
Answer the following four short answer questions in not more than 75 words.\\
निम्नलिखित चार लघु उत्तरीय प्रश्नों के उत्तर अधिकतम 75 शब्दों में दीजिए।
\begin{parts}
  \item Distinguish between old and new concepts of marketing.\\
  विपणन की पुरानी एवं नवीन अवधारणाओं में अन्तर कीजिए।
  \item Classify the products.\\
  उत्पादों का वर्गीकरण कीजिए।
  \item Explain the objectives of pricing.\\
  मूल्यन के उद्देश्यों को समझाइए।
  \item Explain the types of distribution channels.\\
  वितरण वाहिका के प्रकारों को समझाइए।
\end{parts}

\vfill\rule{\linewidth}{0.4pt}
\begin{center}\small
\href{https://bhu-exam.vercel.app/}{Click here to visit our website}\\[4pt]
\href{https://me-aryan.vercel.app/}{Click here to visit our contributor}
\end{center}
\end{document}
